	\documentclass[ ]{article}
\usepackage[ ]{indentfirst}
\title{Revisão P1 \\ Redes de Computadores}
\author{Erickson G. Müller}
\date{27 de Abril de 2026}
\begin{document}
\maketitle

\section*{Conteúdos}
	\begin{itemize}
		\item Redes de Computadores e Internet
		\item Camada de aplicação
		\item Camada de transporte
	\end{itemize}
	
\section{Redes de Computadores}
	O autor adota uma abordagem de camadas \textit{top-down}, com a seguinte estrutura:
	\begin{itemize}
		\item Aplicação
		\item Transporte
		\item Rede
		\item Enlace
		\item Físico
	\end{itemize}
	
	Sinal guiado: com fio
	
	Sinal não-guiado: wireless
	
	Partes de um protocolo(Serviço oferecido; Hipóteses sobre o ambiente onde ele executa; O Vocabulário de mensagens utilizado para implementá-lo; Formato de cada mensagem do vocabulário; Algoritmos)


	\subsection{Bordas da Rede}
	
	\subsection{Comutação de Pacotes}
		A maioria dos comutadores de pacotes utiliza a transmissão \textit{store-and-forward}. O comutador de pacote (\textit{switch}) deve receber o pacote inteiro antes de poder de poder transmitir o primeiro bit para o enlace de saída. Isso acaba gerando uma fila.
	\subsection{Encapsulamento}
		A mensagem (pacote) é enviada na camada de aplicação e é encapsulada para a camada de transporte com o cabeçalho da camada de transporte. Isso é feito para toda a camada até a a camadade enlace. A camada de enlace, com todos os cabeçalhos é enviada até o roteador (?).
\end{document}